% Generated draft: requires mathematical and source review.
\documentclass[11pt]{article}
\usepackage[margin=25mm]{geometry}
\usepackage{fontspec}
\setmainfont{Latin Modern Roman}
\usepackage{amsmath,amssymb,mathtools,unicode-math}
\setmathfont{Latin Modern Math}
\usepackage{enumitem,needspace,xurl,hyperref}
\setlength{\parindent}{0pt}
\setlength{\parskip}{5pt}
\setlength{\emergencystretch}{3em}
\newcommand{\N}{\mathbb{N}}
\newcommand{\Z}{\mathbb{Z}}
\newcommand{\Q}{\mathbb{Q}}
\newcommand{\R}{\mathbb{R}}
\newcommand{\C}{\mathbb{C}}
\newcommand{\F}{\mathbb{F}}
\newcommand{\A}{\mathbb{A}}
\newcommand{\PP}{\mathbb{P}}
\newcommand{\E}{\mathbb{E}}
\newcommand{\eps}{\varepsilon}
\newcommand{\dd}{\,\mathrm{d}}
\begin{document}
\section*{UnsolvedMath ID: 800080}
\textbf{Source problem code:} catalogue entry
\subsection{Definitions and mathematical statement}
A group \(G\) is \emph{simple} if it is nontrivial and its only normal subgroups are \(\{1\}\) and \(G\). A \emph{left-invariant strict total order} on \(G\) is a strict total order \(<\) satisfying
\[
\forall g,x,y\in G,\qquad x<y\Longrightarrow gx<gy.
\]
A group admitting such an order is \emph{left-orderable}. It is \emph{bi-orderable} if it admits a strict total order invariant under both left and right multiplication. The supplied source uses ``orderable'' without specifying which convention is intended.

\textbf{Source question.} Are the Burger--Mozes examples of simple groups orderable?

Under the left-orderability interpretation, for an intended Burger--Mozes simple group \(G\), the property being asked about is
\[
\exists\,<\ \text{a strict total order on }G\quad
\forall g,x,y\in G,\quad x<y\Longrightarrow gx<gy.
\]
The supplied record does not identify the precise examples or construction parameters, nor does it explicitly specify whether this property is requested for every intended example or for at least one of them. Consequently, the displayed property gives a precise conditional interpretation, while the family and its quantification remain unresolved.

The source places this question in the context of finding simple orderable groups of finite cohomological dimension and proposes the Burger--Mozes groups as candidates. For reference, integral cohomological dimension is
\[
\operatorname{cd}_{\Z}(G)
=\operatorname{pd}_{\Z[G]}\Z,
\]
where \(\Z[G]\) is the integral group ring, \(\Z\) is its trivial module, and \(\operatorname{pd}\) denotes projective dimension, with value infinity if no finite projective resolution exists. If integral cohomological dimension and left-orderability are intended, the contextual existence question is
\[
\exists G\quad
G\text{ is simple},\qquad
G\text{ is left-orderable},\qquad
\operatorname{cd}_{\Z}(G)<\infty.
\]
The source does not specify cohomological coefficients. This contextual existence question and the question about the particular Burger--Mozes examples should therefore be retained as distinct targets.
\subsection{Short English statement}
The source asks whether the simple groups constructed by Burger--Mozes are orderable, motivated by the search for simple orderable groups of finite cohomological dimension. The meaning of orderability, the precise examples, the quantification over them, and the cohomological coefficients require clarification. The supplied registry labels the problem open; this draft does not independently verify that status.
\subsection{Sources}
\begin{itemize}
\item Local UnsolvedMath registry: supplied record with permanent ID \(800080\), source problem code \texttt{catalogue entry}, and upstream revision \texttt{d371ee917ac9647fbb06cd6fc529deec6f3dda82}. The record supplies the status ``open'' and identifies Nariman and Yazdi with proposed year 2025.
\item Nariman and Yazdi, \emph{Problem List on Foliations and Diffeomorphism Groups} (2025), Question 21: \url{https://arxiv.org/abs/2503.18274}. The supplied provenance records access on 2026-10-16 and extraction by \texttt{arxiv\_latex}; these are supplied provenance details, not a claim of independent source consultation. The attribution/context is recorded as ``Well-Known,'' so the registry attribution should not be read as establishing original authorship of the question.
\item The source statement cites Burger--Mozes through the reference label \texttt{[burger2000lattices]}. No corresponding full bibliographic entry or link is supplied, so the exact referenced construction cannot be identified from this record alone.
\end{itemize}
\end{document}
